% =============================================================================
% hyphenation-id.tex
%
% Aturan pemenggalan kata (hyphenation) Bahasa Indonesia untuk kata-kata yang
% sering muncul di skripsi IPB dan tidak ditangani baik oleh babel default.
%
% Paket: hyphenat (\hyphenation{...}). Tambahkan kata di dalam argumen,
%        pisahkan dengan spasi, gunakan tanda "-" di lokasi pemenggalan.
%
% Catatan: Daftar ini hanya pelengkap; sebagian besar kata sudah ditangani
%          oleh babel[indonesian]. Tambahkan kata yang menurut Anda sering
%          dipenggal salah di sini.
% =============================================================================

\hyphenation{
  agri-kultur
  agro-nomi
  bio-teknolo-gi
  eko-sis-tem
  eko-no-mi
  fito-plank-ton
  herbari-um
  in-sti-tut
  kon-servasi
  ling-kung-an
  manaje-men
  mikro-or-ga-nis-me
  mo-le-kul
  neu-ron
  nu-tri-si
  or-ga-nis-me
  pa-ra-sit
  per-tanian
  per-ubah-an
  pol-u-si
  pop-u-lasi
  sum-ber-daya
  ta-nah
  ta-na-man
  tek-nolo-gi
  varia-bel
  vi-rologi
  zo-oplank-ton
}
